\section{Object-Oriented Programming Language}

The project uses Java as its object-oriented programming language.
The source code targets Java 21, and the build sets this version through the Maven compiler release option.
Java provides classes, interfaces, abstract classes, inheritance, and polymorphism, which the design of the simulator relies on.
Table~\ref{tab:oo} links these language features to the parts of the system that use them.

\begin{table}[H]
\centering
\small
\caption{Java language features and their use in the system}
\label{tab:oo}
\begin{tabular}{>{\raggedright\arraybackslash}p{4.2cm}>{\raggedright\arraybackslash}p{10.8cm}}
\toprule
\textbf{Language feature} & \textbf{Use in the system} \\
\midrule
Abstract classes & \texttt{DigitalComponent} and \texttt{LogicGate} define the shared behavior of all components and gates. \\
Inheritance & Each gate, such as \texttt{AndGate}, extends \texttt{LogicGate} and implements only its own Boolean function. \\
Interfaces & \texttt{SignalObserver}, \texttt{MemoryState}, and \texttt{ComponentCreator} support the Observer, State, and Factory Method patterns. \\
Polymorphism & The engine calls \texttt{evaluate} on any component without knowledge of its concrete class. \\
Enumerations & \texttt{VoltageState} represents the four signal values 0, 1, X, and Z. \\
Encapsulation & Each \texttt{Wire} hides its list of drivers and observers behind a small set of methods. \\
\bottomrule
\end{tabular}
\end{table}

The desktop interface uses Java Swing, which is part of the standard Java platform.
The browser version runs the same Java engine after the TeaVM compiler translates it to JavaScript.
A single HTML page, \texttt{index.html}, provides the browser interface.
The project writes no simulation logic in JavaScript.
